\documentclass[10pt,letterpaper]{article}
\usepackage[T1]{fontenc}\usepackage[utf8]{inputenc}\usepackage[scaled=0.95]{helvet}\renewcommand{\familydefault}{\sfdefault}
\usepackage[letterpaper,top=0.55in,bottom=0.6in,left=0.7in,right=0.7in]{geometry}
\usepackage{microtype}\usepackage{enumitem}\usepackage{xcolor}\usepackage{fancyhdr}\usepackage[hidelinks]{hyperref}
\newcommand{\docversion}{2026-10-23-10-58-02} % set by build.py (America/Toronto)
\hypersetup{pdftitle={St Monica's parish finances and the Terrebonne bike lane: executive summary},pdfsubject={version \docversion}}
\setlength{\parindent}{0pt}\setlength{\parskip}{0.25em}
\pagestyle{fancy}\fancyhf{}\renewcommand{\headrulewidth}{0pt}
\fancyfoot[C]{\scriptsize\color{gray}A resident of Côte-des-Neiges--Notre-Dame-de-Grâce, from public records. Full note, data and scripts: github.com/OpenMobilityData/Terrebonne-open-sources, folder 2024-st-monica-cra-t3010. {\tiny v.\,\docversion}}
\newcommand{\B}[1]{\textbf{#1}}
\begin{document}
{\LARGE\bfseries St Monica's parish finances and the Terrebonne bike lane}\par\smallskip
{\large Executive summary of the data note, 2008--2024}\par\medskip

\textbf{Question.} Did the protected bike lane built on rue de Terrebonne in May--June 2024 have a measurable effect on the parish's finances? \textbf{Record.} The public financial section of the parish's \B{17} charity returns (CRA form T3010, 2008--2024), the parish's own statements published in its bulletins (2017--March 2023), all \B{519} bulletins of its website (2014--October 18, 2026), the returns of the \B{12} other Catholic parishes of the same postal areas, and the parish's Facebook page. Fiscal years are calendar years; the 2024 return is the only one with the lane, for seven months; the 2025 return is not yet published.

\textbf{What the 2024 return shows}
\begin{itemize}[nosep,leftmargin=1.2em]
\item Total revenue \B{\$250,718}, \B{6~\%} below 2023 and \B{4~\%} below 2022, above 2020 and 2021; within the 2019--2023 range (\$202,900 to \$276,981).
\item Receipted gifts \B{\$201,562} (\B{$-$2~\%} on 2023); gifts without receipt \B{\$18,311}; rental income \B{\$15,339} (level with 2022--2023); « other revenue » \B{\$15,506}, described by the parish as faith-education fees and vigil lights, \emph{not} funerals or weddings.
\item Expenditures \B{\$269,373}; the gap of \B{\$18,655} is the third deficit in a row after 2022 (\B{\$31,771}) and 2023 (\B{\$8,240}), and the seventh of 17 years. Its increase over 2023 (\$10,415) is smaller than the rise in compensation (\B{\$15,783}).
\item Cash \B{\$197,111} (from \$31,528) and liabilities \B{\$180,109} (from \$2,123): the highest in the series, unexplained by the return.
\end{itemize}

\textbf{Weddings and funerals}
\begin{itemize}[nosep,leftmargin=1.2em]
\item The return has no line for them. The line the parish once labelled « masses, funeral, weddings » carried \B{\$32,913 to \$63,217} a year in 2008--2014 and has been nearly empty since 2015, the year a line of unreceipted gifts appeared. The parish's own statements group hall rentals, bequests, funerals and weddings in one « other income » line (\B{\$83,157} in 2019). Such sums can sit on four return lines; nothing separates them, and no count of weddings or funerals exists in any public source.
\end{itemize}

\textbf{Against the twelve neighbouring parishes (gifts = both gift lines; each parish indexed to its own 2017--2019 mean)}
\begin{itemize}[nosep,leftmargin=1.2em]
\item 2023 to 2024: the parish's gifts \B{$-$6.1~\%} vs the peers' median \B{$+$0.8~\%} (rank \B{10 of 13}; seven peers up, five down, three by more); total revenue \B{$-$6.0~\%} vs \B{$+$5.0~\%} (rank 9). Both inside the middle half of the peers' changes, whose full range runs from $-$50 to $+$73~\%.
\item The parish was above the peers' median in 2022 and below from 2023: the gap opened \emph{before} the lane.
\end{itemize}

\textbf{What this record could detect}
\begin{itemize}[nosep,leftmargin=1.2em]
\item Pre-lane year-to-year variation (2015--2023) sets a threshold of about \B{15 to 26~\%} of a year's gifts (parish alone) or \B{14 to 16 points} against the peers' median. With seven months of lane in 2024, a lane effect would have had to remove \B{25 to 45~\%} of the parish's full-year gifts to stand apart. Observed: $-$6~\%. The 2025 return lowers the threshold to \B{14 to 26~\%}; two full years to \B{10 to 18~\%}.
\end{itemize}

\textbf{What the parish itself has said}
\begin{itemize}[nosep,leftmargin=1.2em]
\item Its statements and the return agree to the dollar for 2019, 2020 and 2022; the monthly statements stopped in \B{March 2023}, fourteen months before the lane, and never resumed.
\item In the \B{117} bulletins published since July 2024: \B{0} references to the lane, \B{0} to financial difficulties linked to it, \B{0} to parking difficulties. The same bulletins carry a 75th-anniversary repairs campaign (\B{\$163,858} raised in 2024--25, a \$25,000 matching pledge in 2026), civic notices (a municipal-election appeal, petitions, legislation), the MacKay Center weekend-parking passes (announced \B{October 15, 2023}, seven months before the lane), and a « Save St.~Monica's » rally on \B{November 24, 2024} announced without any stated reason and without any published result.
\item The Facebook page campaigned against the 2020 temporary lane (petition, council questions); since the 2024 lane its one official comment on bike lanes (August 14, 2024) shared an article in which Philadelphia congregations gave up bike-lane parking to stop alienating their neighbours. It has never stated a financial effect.
\end{itemize}

\textbf{Conclusion.} No figure in the 2024 return stands apart from the parish's own year-to-year variation or from its neighbours'; the deficits precede the lane; the lines nearest the claims are within or above their recent range. A measurable impact is \B{neither established nor excluded}: an effect would have had to be of the order of a quarter of a year's gifts to show in a seven-month return. The test rests on the 2025 return, read the same way against the peers, and on monthly figures the parish no longer publishes.
\end{document}
